\documentclass[aspectratio=169,10pt]{beamer}

% ---------- Theme ----------
\usepackage{beamerthemeAcademic}
\useblue

% ---------- Fonts ----------
\usepackage{fontspec}
\usepackage{xeCJK}
\IfFontExistsTF{Microsoft YaHei}{
  \setCJKmainfont{Microsoft YaHei}
  \setCJKsansfont{Microsoft YaHei}
  \setCJKmonofont{Microsoft YaHei}
}{
  \IfFontExistsTF{SimSun}{
    \setCJKmainfont{SimSun}
    \setCJKsansfont{SimSun}
    \setCJKmonofont{SimSun}
  }{
    \setCJKmainfont{Noto Serif CJK SC}
    \setCJKsansfont{Noto Sans CJK SC}
    \setCJKmonofont{Noto Sans CJK SC}
  }
}
\IfFontExistsTF{Times New Roman}{\setmainfont{Times New Roman}}{\setmainfont{TeX Gyre Termes}}
\IfFontExistsTF{Arial}{\setsansfont{Arial}}{\setsansfont{TeX Gyre Heros}}
\IfFontExistsTF{Consolas}{\setmonofont{Consolas}}{\setmonofont{TeX Gyre Cursor}}

% ---------- Packages ----------
\usepackage{amsmath,amssymb}
\usepackage{booktabs}
\usepackage{colortbl}
\usepackage{array}
\usepackage{tabularx}
\usepackage{makecell}
\usepackage{pifont}
\usepackage{tikz}
\usetikzlibrary{arrows.meta,positioning,calc,fit,shapes.geometric}

\setlength{\emergencystretch}{2em}
\hfuzz=3pt
\setlength{\parskip}{2pt}
\AtBeginSection[]{}

% ---------- Local styles ----------
\definecolor{rvgreen}{RGB}{24,118,67}
\definecolor{rvred}{RGB}{164,42,42}
\definecolor{rvorange}{RGB}{181,103,22}
\definecolor{rvblue2}{RGB}{41,91,140}
\newcommand{\cmark}{\textcolor{rvgreen}{\ding{51}}}
\newcommand{\warnmark}{\textcolor{rvorange}{\ding{115}}}
\newcommand{\xmark}{\textcolor{rvred}{\ding{55}}}
\newcommand{\ok}[1]{\textcolor{rvgreen}{\textbf{#1}}}
\newcommand{\warn}[1]{\textcolor{rvorange}{\textbf{#1}}}
\newcommand{\nope}[1]{\textcolor{rvred}{\textbf{#1}}}
\newcommand{\metric}[2]{\begin{tabular}{@{}l@{}}{\Large\bfseries\color{accentcolor}#1}\\[-2pt]{\scriptsize #2}\end{tabular}}
\newcolumntype{Y}{>{\raggedright\arraybackslash}X}
\newcolumntype{L}[1]{>{\raggedright\arraybackslash}p{#1}}

\title[RV-MalTrace 阶段性进展]{RV-MalTrace 阶段性进展}
\author[RV-MalTrace]{RV-MalTrace}
\institute[阶段性汇报]{阶段性汇报}
\date{2026年6月12日}
\setsupervisor{}
\setmajor{}

\begin{document}

% 1
\begin{frame}[plain]
  \begin{tikzpicture}[remember picture,overlay]
    \fill[accentcolor] (current page.north west) rectangle ([yshift=-0.16cm]current page.north east);
    \fill[accentcolor!7] ([yshift=-0.16cm]current page.north west) rectangle (current page.south east);
    \node[anchor=west,text=accentcolor,font=\sffamily\bfseries\fontsize{25pt}{32pt}\selectfont]
      at ([xshift=1.0cm,yshift=-1.55cm]current page.north west)
      {RV-MalTrace 阶段性进展};
    \node[anchor=west,text=inkcolor,font=\sffamily\fontsize{17pt}{24pt}\selectfont]
      at ([xshift=1.0cm,yshift=-2.45cm]current page.north west)
      {Genesys2/CVA6 板级记录、行为解释与论文推进};
    \draw[accentcolor,line width=0.8pt]
      ([xshift=1.0cm,yshift=-3.15cm]current page.north west) --
      ([xshift=-1.0cm,yshift=-3.15cm]current page.north east);
    \node[anchor=west,align=left,text=inkcolor,text width=13.4cm,font=\sffamily\fontsize{11pt}{16pt}\selectfont]
      at ([xshift=1.05cm,yshift=-4.35cm]current page.north west) {
        本次重点：整体框架、当前进度、成果数据、结论边界、下一步安排\\[4pt]
        口径说明：只讲已验证的安全样例与捕获窗口；\\
        不把结果扩展为真实恶意软件检测准确率
      };
    \node[anchor=south east,text=textgray,font=\sffamily\fontsize{9pt}{12pt}\selectfont]
      at ([xshift=-1.0cm,yshift=0.75cm]current page.south east) {2026年6月12日};
  \end{tikzpicture}
\end{frame}

% 2
\begin{frame}{汇报提纲}
  \vskip0.15cm
  \begin{center}
  \renewcommand{\arraystretch}{1.28}
  \begin{tabularx}{0.92\linewidth}{@{}L{0.9cm}Y@{}}
    \toprule
    \textbf{\color{accentcolor}一} & \textbf{整体框架}：系统做什么，trace 怎样变成可复查结果。 \\
    \textbf{\color{accentcolor}二} & \textbf{当前进度}：从 35T/VexRiscv 小容量结果推进到 Genesys2/CVA6。 \\
    \textbf{\color{accentcolor}三} & \textbf{成果数据}：P0、BRAM、safe surrogate、评估矩阵、资源与运行时间。 \\
    \textbf{\color{accentcolor}四} & \textbf{结论边界}：现在可以说什么，哪些还不能说。 \\
    \textbf{\color{accentcolor}五} & \textbf{下一阶段}：补哪些实验，论文主线怎样收束。 \\
    \bottomrule
  \end{tabularx}
  \end{center}
  \vskip0.2cm
  \keybox{\textbf{汇报原则：}每个结论都对应检查文件或结果摘要；不使用没有数据支撑的说法。}
\end{frame}

\section{整体框架}
\sectiondivider{一}{整体框架}

% 4
\begin{frame}{这个系统到底在做什么}
  \vskip0.1cm
  RV-MalTrace 的研究对象不是一个完整的“检测器”，而是一个\alert{硬件辅助行为记录原型}。它在处理器提交路径旁边记录关键事件，再在主机侧把记录解释成系统调用、路径、进程、动态映射等行为证据。

  \vskip0.25cm
  \begin{center}
  \renewcommand{\arraystretch}{1.20}
  \begin{tabularx}{0.92\linewidth}{@{}L{2.8cm}Y@{}}
    \toprule
    \textbf{输入} & CVA6/RISC-V 的提交、异常、特权级、系统调用相关信号。 \\
    \rowcolor{black!4}
    \textbf{硬件侧输出} & 固定格式 trace 记录、参数快照、窗口标记、DROP 计数。 \\
    \textbf{主机侧输出} & 解码后的 JSONL、行为图、评估表和不能越界的结论说明。 \\
    \rowcolor{black!4}
    \textbf{设计原则} & trace 满时丢记录并计数，不停住处理器执行。 \\
    \bottomrule
  \end{tabularx}
  \end{center}
\end{frame}

% 5
\begin{frame}{从处理器到结果文件}
  \vskip0.05cm
  \begin{center}
  \begin{tikzpicture}[
    scale=0.92,transform shape,
    node distance=0.45cm and 0.72cm,
    box/.style={draw=lightline,fill=white,rounded corners=2pt,minimum height=0.74cm,align=center,font=\scriptsize,text width=2.38cm},
    wide/.style={draw=accentcolor!60,fill=accentcolor!8,rounded corners=2pt,minimum height=0.78cm,align=center,font=\scriptsize,text width=2.75cm},
    arrow/.style={-{Latex[length=2.0mm]},line width=0.55pt,color=accentcolor}
  ]
    \node[wide] (core) {CVA6\\RISC-V core};
    \node[box,right=of core] (tap) {提交 / 异常 / 系统调用\\旁路采样};
    \node[box,right=of tap] (packet) {事件打包\\过滤与编号};
    \node[box,right=of packet] (sink) {BRAM / ILA\\DROP 计数};

    \node[box,below=0.95cm of tap] (decode) {主机解码\\JSONL};
    \node[box,right=of decode] (align) {进程归属\\对照日志对齐};
    \node[box,right=of align] (sem) {行为解释\\路径 / 进程 / mmap};
    \node[wide,below=0.95cm of align] (report) {评估矩阵\\结论边界};

    \draw[arrow] (core) -- (tap);
    \draw[arrow] (tap) -- (packet);
    \draw[arrow] (packet) -- (sink);
    \draw[arrow] (sink) |- (decode);
    \draw[arrow] (decode) -- (align);
    \draw[arrow] (align) -- (sem);
    \draw[arrow] (sem) |- (report);
  \end{tikzpicture}
  \figcap{图 1\;系统把硬件记录和离线解释分开，硬件负责事实记录，主机侧负责语义恢复与检查。}
  \end{center}
  \vskip0.08cm
  \keybox{\textbf{关键点：}当前能汇报的是“哪些行为被记录并可复查”，不是“能检测所有恶意软件”。}
\end{frame}

% 6
\begin{frame}{一条 trace 怎样变成结论}
  \vskip0.05cm
  \begin{center}
  \renewcommand{\arraystretch}{1.22}
  \begin{tabularx}{0.94\linewidth}{@{}L{1.4cm}L{2.4cm}Y@{}}
    \toprule
    \textbf{步骤} & \textbf{产物} & \textbf{检查点} \\
    \midrule
    01 & BRAM/ILA 窗口 & marker 是否存在，窗口是否属于目标样例。 \\
    \rowcolor{black!4}
    02 & JSONL 事件 & 事件能否 parse，序号是否连续，syscall 是否成对。 \\
    03 & 运行时归属 & trace 是否能归到目标进程和目标 ELF。 \\
    \rowcolor{black!4}
    04 & 行为解释 & 路径、进程、动态映射、反调试节点是否能生成。 \\
    05 & 结果摘要 & PASS、DROP、wrap、不能声称项是否记录清楚。 \\
    \bottomrule
  \end{tabularx}
  \end{center}
  \vskip0.15cm
  \keybox{这条流程的价值在于可以回到文件复查，不依赖人工说“看起来是对的”。}
\end{frame}

% 7
\begin{frame}{记录满了怎么办}
  \vskip0.1cm
  \begin{columns}[T,onlytextwidth]
    \column{0.48\textwidth}
    \begin{block}{不采用的做法}
      停住处理器等 trace 缓冲区。这样会改变被观察程序的执行节奏，也会让安全分析结果很难解释。
    \end{block}
    \column{0.48\textwidth}
    \begin{block}{当前采用的做法}
      trace 侧记录 DROP，处理器继续执行。汇报时只解释捕获窗口内已记录且未丢失的结果。
    \end{block}
  \end{columns}
  \vskip0.25cm
  \begin{center}
  \renewcommand{\arraystretch}{1.18}
  \begin{tabularx}{0.88\linewidth}{@{}L{2.8cm}Y@{}}
    \toprule
    \textbf{好处} & 不把监测机制变成程序运行的一部分，便于说明低扰动。 \\
    \rowcolor{black!4}
    \textbf{代价} & 结论必须限定在捕获窗口，不能声称全程无遗漏。 \\
    \bottomrule
  \end{tabularx}
  \end{center}
\end{frame}

% 8
\begin{frame}{结果不是一句“通过”}
  \vskip0.05cm
  \begin{center}
  \renewcommand{\arraystretch}{1.18}
  \begin{tabularx}{0.94\linewidth}{@{}L{2.4cm}Y@{}}
    \toprule
    \textbf{结果文件} & \textbf{能复查的内容} \\
    \midrule
    trace 记录 & 每个样例的窗口内事件、序号、系统调用、DROP、wrap。 \\
    \rowcolor{black!4}
    运行时归属 & 目标进程路径、目标 ELF、子进程归属。 \\
    行为图 & 路径访问、进程创建、动态映射、反调试行为节点。 \\
    \rowcolor{black!4}
    对照结果 & strace、QEMU、软件插桩、event-only、pointer snapshot 等对照。 \\
    结论边界 & 能说什么、不能说什么，以及为什么。 \\
    \bottomrule
  \end{tabularx}
  \end{center}
  \vskip0.1cm
  \keybox{当前工作的重点已经从“能不能跑”转到“结果能不能被别人复查”。}
\end{frame}

\section{当前进度}
\sectiondivider{二}{当前进度}

% 10
\begin{frame}{这周最大的变化}
  \vskip0.1cm
  5 月底的主要结果是 35T/VexRiscv 小容量矩阵。6 月 11 日之后，汇报重点已经转到 \alert{Genesys2/CVA6}：P0 样例、BRAM 窗口重复、safe surrogate、行为解释和对照评估都进入当前结果目录。

  \vskip0.25cm
  \begin{center}
  \renewcommand{\arraystretch}{1.18}
  \begin{tabularx}{0.94\linewidth}{@{}L{2.8cm}Y@{}}
    \toprule
    \textbf{2026-05-21} & 35T 小容量矩阵 13/13 通过，证明受控样例能在 512 条记录预算内完成检查。 \\
    \rowcolor{black!4}
    \textbf{2026-05-29} & 周报版证据链成型，开始整理论文可用结果和边界。 \\
    \textbf{2026-06-11} & Genesys2/CVA6 P0 结果通过，包含系统调用配对和运行时归属。 \\
    \rowcolor{black!4}
    \textbf{2026-06-12} & P0 40 次、safe surrogate 80 次 BRAM 窗口重复进入当前汇报。 \\
    \bottomrule
  \end{tabularx}
  \end{center}
\end{frame}

% 11
\begin{frame}{当前进度分三层看}
  \vskip0.1cm
  \begin{columns}[T,onlytextwidth]
    \column{0.32\textwidth}
    \begin{block}{板级采集}
      Genesys2/CVA6 能在 marker 窗口里采到连续事件，P0 和安全替代样例都有重复结果。
    \end{block}
    \column{0.32\textwidth}
    \begin{block}{结果复查}
      系统调用配对、DROP、wrap、序号连续性都有检查，并写入摘要文件。
    \end{block}
    \column{0.32\textwidth}
    \begin{block}{行为解释}
      路径、进程、动态映射、反调试节点和对照方法进入同一个评估表。
    \end{block}
  \end{columns}
  \vskip0.25cm
  \keybox{目前最重要的进展是：CVA6 板级结果不再只是“抓到事件”，而是已经连接到可解释的结果数据。}
\end{frame}

% 12
\begin{frame}{P0 板级连续记录结果}
  \vskip0.05cm
  \begin{center}
  \renewcommand{\arraystretch}{1.20}
  \begin{tabularx}{0.92\linewidth}{@{}L{3.0cm}L{2.3cm}Y@{}}
    \toprule
    \textbf{检查项} & \textbf{结果} & \textbf{解释} \\
    \midrule
    P0 样例 & \ok{4/4} & \texttt{hello\_write}、\texttt{file\_open\_read\_write}、\texttt{fork\_exec}、\texttt{illegal\_instruction}。 \\
    \rowcolor{black!4}
    未解释 DROP & \ok{0} & 捕获窗口内没有无法解释的 DROP。 \\
    运行时归属 & \ok{YES} & 结果能归属到板上运行的目标进程和 ELF。 \\
    \rowcolor{black!4}
    严格检查 & \ok{PASS} & syscall entry/return 使用非零 \texttt{syscall\_id} 配对。 \\
    \bottomrule
  \end{tabularx}
  \end{center}
  \vskip0.12cm
  \keybox{这部分支撑的说法是：Genesys2/CVA6 上 P0 安全样例的窗口内记录已经可复查。}
\end{frame}

% 13
\begin{frame}{BRAM 窗口重复结果}
  \vskip0.05cm
  \begin{center}
  \begin{tabular}{ccc}
    \metric{40}{P0 重复窗口} &
    \metric{80}{safe surrogate 重复窗口} &
    \metric{120}{合计接受窗口}
  \end{tabular}
  \end{center}
  \vskip0.25cm
  \begin{center}
  \renewcommand{\arraystretch}{1.20}
  \begin{tabularx}{0.90\linewidth}{@{}L{3.3cm}Y@{}}
    \toprule
    \textbf{P0 BRAM 摘要} & 4 个 P0 样例，每个 10 次；总尝试 41 次，接受 40 次；最大未解释 DROP 为 0，最大 wrap 为 0。 \\
    \rowcolor{black!4}
    \textbf{安全替代样例摘要} & 8 个样例，每个 10 次；总计 80 次；最大未解释 DROP 为 0，最大 wrap 为 0。 \\
    \bottomrule
  \end{tabularx}
  \end{center}
  \vskip0.1cm
  \keybox{这些数字只说明接受窗口的稳定性，不代表系统全程没有任何遗漏。}
\end{frame}

% 14
\begin{frame}{safe surrogate 8 个样例}
  \vskip0.05cm
  \begin{center}\scriptsize
  \renewcommand{\arraystretch}{1.15}
  \begin{tabularx}{0.96\linewidth}{@{}L{3.4cm}Y L{2.0cm}@{}}
    \toprule
    \textbf{样例} & \textbf{覆盖的行为} & \textbf{BRAM 重复} \\
    \midrule
    \texttt{file\_scan} & 目录打开、读取、关闭。 & 10/10 \\
    \rowcolor{black!4}
    \texttt{batch\_open\_read\_write} & 批量文件读写。 & 10/10 \\
    \texttt{self\_copy\_sim} & 自复制式文件流。 & 10/10 \\
    \rowcolor{black!4}
    \texttt{abnormal\_syscall\_sequence} & 异常系统调用序列。 & 10/10 \\
    \texttt{illegal\_trap} & 非法指令触发异常。 & 10/10 \\
    \rowcolor{black!4}
    \texttt{process\_chain} & fork/exec/wait 进程链。 & 10/10 \\
    \texttt{dynamic\_executable\_memory} & 动态可执行映射。 & 10/10 \\
    \rowcolor{black!4}
    \texttt{anti\_debug\_like} & ptrace 风格的反调试可见性。 & 10/10 \\
    \bottomrule
  \end{tabularx}
  \end{center}
  \vskip0.08cm
  \keybox{这里的样例是安全替代样例，不是真实恶意载荷。}
\end{frame}

% 15
\begin{frame}{12 个样例进入同一评估表}
  \vskip0.05cm
  \begin{center}
  \begin{tabular}{ccc}
    \metric{12}{评估样例} &
    \metric{6}{对照方法} &
    \metric{PASS}{矩阵状态}
  \end{tabular}
  \end{center}
  \vskip0.22cm
  \begin{center}\scriptsize
  \renewcommand{\arraystretch}{1.15}
  \begin{tabularx}{0.94\linewidth}{@{}L{4.0cm}Y@{}}
    \toprule
    \textbf{对照方法} & \textbf{用途} \\
    \midrule
    \texttt{rv\_maltrace\_event\_only} & 只用事件记录，作为硬件记录的基础版本。 \\
    \rowcolor{black!4}
    \texttt{rv\_maltrace\_pointer\_snapshot} & 有界 pointer 前缀记录，检查哪些参数能被硬件补充。 \\
    \texttt{rv\_maltrace\_kernel\_helper} & 受信内核 helper，作为语义补充路线。 \\
    \rowcolor{black!4}
    \texttt{strace} / \texttt{qemu\_strace} & 软件 syscall 对照。 \\
    \texttt{software\_instrumentation} & 编译或插桩路线对照。 \\
    \bottomrule
  \end{tabularx}
  \end{center}
\end{frame}

\section{成果数据}
\sectiondivider{三}{成果数据}

% 17
\begin{frame}{行为恢复指标}
  \vskip0.05cm
  \begin{center}\small
  \renewcommand{\arraystretch}{1.18}
  \begin{tabularx}{0.92\linewidth}{@{}L{4.2cm}L{2.0cm}Y@{}}
    \toprule
    \textbf{指标} & \textbf{结果} & \textbf{含义} \\
    \midrule
    expected syscall recall & 1.0 & 预期系统调用在受控样例中被恢复。 \\
    \rowcolor{black!4}
    syscall precision & 1.0 & 当前恢复出的系统调用没有超出预期集合。 \\
    argument reconstruction accuracy & 1.0 & 受控样例中的参数恢复与对照一致。 \\
    \rowcolor{black!4}
    behavior rule recall & 1.0 & 行为规则在当前样例中被触发。 \\
    anti-analysis visibility & 1.0 & 反调试相关行为在样例中可见。 \\
    \rowcolor{black!4}
    benign false positive rate & 0.0 & 当前安全样例没有被误报为恶意。 \\
    unaccounted drop & 0 & 当前摘要中的窗口没有未解释 DROP。 \\
    \bottomrule
  \end{tabularx}
  \end{center}
  \vskip0.1cm
  \keybox{这些指标只适用于当前受控样例集，不能外推为真实恶意软件检测准确率。}
\end{frame}

% 18
\begin{frame}{语义恢复的来源要讲清楚}
  \vskip0.1cm
  \begin{columns}[T,onlytextwidth]
    \column{0.48\textwidth}
    \begin{block}{现在可以说}
      路径、进程、动态映射等行为能在当前样例中恢复，并写入行为图和语义事件文件。
    \end{block}
    \column{0.48\textwidth}
    \begin{block}{需要说明}
      完整字符串语义主要来自受信 QEMU/strace companion 的对齐；硬件 pointer snapshot 是有界前缀信息，不是完整字符串导出。
    \end{block}
  \end{columns}
  \vskip0.25cm
  \begin{center}
  \renewcommand{\arraystretch}{1.18}
  \begin{tabularx}{0.90\linewidth}{@{}L{3.0cm}Y@{}}
    \toprule
    \textbf{硬件记录} & syscall、返回、异常、上下文、少量参数快照、DROP。 \\
    \rowcolor{black!4}
    \textbf{对齐补充} & 路径字符串、进程上下文、对照日志中的语义信息。 \\
    \bottomrule
  \end{tabularx}
  \end{center}
\end{frame}

% 19
\begin{frame}{代表性行为图结果}
  \vskip0.05cm
  \begin{center}\scriptsize
  \renewcommand{\arraystretch}{1.16}
  \begin{tabularx}{0.96\linewidth}{@{}L{3.3cm}Y L{2.2cm}@{}}
    \toprule
    \textbf{样例} & \textbf{恢复出的行为} & \textbf{状态} \\
    \midrule
    \texttt{file\_scan} & 路径节点、目录读取节点、close 节点可复查。 & \ok{PASS} \\
    \rowcolor{black!4}
    \texttt{process\_chain} & fork/exec/wait 关系可见，子进程归属可证明。 & \ok{PASS} \\
    动态可执行映射 & mmap/mprotect 风格的动态可执行映射节点可见。 & \ok{PASS} \\
    \rowcolor{black!4}
    \texttt{anti\_debug\_like} & ptrace 风格反调试行为在对照中可见。 & \ok{PASS} \\
    异常 syscall 序列 & 故意失败的 fd 被标成预期异常行为。 & \ok{PASS} \\
    \bottomrule
  \end{tabularx}
  \end{center}
  \vskip0.1cm
  \keybox{这些是“行为解释结果”，不是恶意家族分类结果。}
\end{frame}

% 20
\begin{frame}{进程和 ELF 归属}
  \vskip0.1cm
  \begin{center}
  \begin{tabular}{ccc}
    \metric{12/12}{目标进程归属} &
    \metric{2}{子进程样例} &
    \metric{separate}{动态库区分}
  \end{tabular}
  \end{center}
  \vskip0.25cm
  \begin{center}
  \renewcommand{\arraystretch}{1.20}
  \begin{tabularx}{0.92\linewidth}{@{}L{3.0cm}Y@{}}
    \toprule
    \textbf{目标进程} & 每个样例都有目标进程路径，trace 不只是裸事件。 \\
    \rowcolor{black!4}
    \textbf{子进程} & \texttt{fork\_exec} 和 \texttt{process\_chain} 能证明子进程目标。 \\
    \textbf{动态库} & 动态库事件不会被直接误归到目标 ELF。 \\
    \bottomrule
  \end{tabularx}
  \end{center}
\end{frame}

% 21
\begin{frame}{代码归属要分两层讲}
  \vskip0.1cm
  \begin{center}
  \renewcommand{\arraystretch}{1.24}
  \begin{tabularx}{0.92\linewidth}{@{}L{3.0cm}Y@{}}
    \toprule
    \textbf{可以说} & 当前 board trace 通过 ELF symbol/range 和运行时进程图做到函数级归属，相关摘要中的 function-level attribution 可用。 \\
    \rowcolor{black!4}
    \textbf{不能说} & 当前板上 trace 已经直接带出 DWARF 源代码行号。source-line 结果来自 debug/no-PIE 的 source-equivalent sidecar。 \\
    \bottomrule
  \end{tabularx}
  \end{center}
  \vskip0.15cm
  \keybox{这页建议主动讲，避免老师追问时把函数级归属和源码行号混在一起。}
\end{frame}

% 22
\begin{frame}{资源与运行时间}
  \vskip0.05cm
  \begin{center}\scriptsize
  \renewcommand{\arraystretch}{1.16}
  \begin{tabular}{@{}lrrr@{}}
    \toprule
    \textbf{指标} & \textbf{Baseline} & \textbf{Trace-enabled} & \textbf{Delta} \\
    \midrule
    LUT & 84,923 & 106,428 & +21,505 (+25.32\%) \\
    \rowcolor{black!4}
    FF & 56,491 & 65,634 & +9,143 (+16.18\%) \\
    BRAM18 equiv & 108 & 114 & +6 (+5.56\%) \\
    \rowcolor{black!4}
    DSP & 27 & 27 & 0 \\
    Slack & 0.177 ns & 0.177 ns & 0 \\
    \bottomrule
  \end{tabular}
  \end{center}
  \vskip0.2cm
  \begin{center}
  \renewcommand{\arraystretch}{1.15}
  \begin{tabularx}{0.90\linewidth}{@{}L{3.2cm}Y@{}}
    \toprule
    \textbf{运行时间统计} & BRAM、event-only、pointer-disabled 三种 production 模式相对 trace-off 的 median slowdown 都是 1.0x。 \\
    \rowcolor{black!4}
    \textbf{解释边界} & 这是当前 UART START/DONE marker 的 wall-clock 统计，不应写成所有 workload 的性能结论。 \\
    \bottomrule
  \end{tabularx}
  \end{center}
\end{frame}

% 23
\begin{frame}{DROP 与 wrap 结果}
  \vskip0.1cm
  \begin{center}
  \begin{tabular}{ccc}
    \metric{0}{未解释 DROP} &
    \metric{0}{wrap count} &
    \metric{PASS}{窗口检查}
  \end{tabular}
  \end{center}
  \vskip0.25cm
  \begin{center}
  \renewcommand{\arraystretch}{1.20}
  \begin{tabularx}{0.90\linewidth}{@{}L{3.0cm}Y@{}}
    \toprule
    \textbf{当前结果} & 所有列入当前摘要的 Genesys2/CVA6 捕获窗口，未解释 DROP 和 wrap 都为 0。 \\
    \rowcolor{black!4}
    \textbf{正确表达} & 这是捕获窗口内的正确性和完整性检查，不是全系统连续语义恢复。 \\
    \bottomrule
  \end{tabularx}
  \end{center}
\end{frame}

\section{结论边界}
\sectiondivider{四}{结论边界}

% 25
\begin{frame}{现在可以稳妥讲的结论}
  \vskip0.05cm
  \begin{center}\scriptsize
  \renewcommand{\arraystretch}{1.18}
  \begin{tabularx}{0.96\linewidth}{@{}L{3.4cm}Y@{}}
    \toprule
    \textbf{结论} & \textbf{精确含义} \\
    \midrule
    CVA6 板级证据 & 安全 P0 与 safe surrogate 样例在 Genesys2/CVA6 上有窗口内记录。 \\
    \rowcolor{black!4}
    BRAM 窗口稳定 & 120 个接受窗口具备序号连续、无 wrap、无未解释 DROP。 \\
    行为可解释 & 路径、进程、动态映射、反调试节点能生成并复查。 \\
    \rowcolor{black!4}
    对照可复查 & strace、QEMU、软件插桩、event-only、pointer snapshot 等方法进入同一评估表。 \\
    论文材料可整理 & 结果、通过条件、不能声称项已经能组成阶段性报告。 \\
    \bottomrule
  \end{tabularx}
  \end{center}
\end{frame}

% 26
\begin{frame}{不能把结果说大}
  \vskip0.05cm
  \begin{center}\scriptsize
  \renewcommand{\arraystretch}{1.18}
  \begin{tabularx}{0.96\linewidth}{@{}L{3.6cm}L{1.5cm}Y@{}}
    \toprule
    \textbf{不能声称} & \textbf{状态} & \textbf{原因} \\
    \midrule
    真实恶意软件检测准确率 & \nope{NO} & 当前样例是安全合成样例或 syscall-only surrogate，不是真实恶意载荷。 \\
    \rowcolor{black!4}
    家族分类或 IOC 覆盖 & \nope{NO} & 行为指标说明受控工作负载内的恢复结果，不代表家族识别质量。 \\
    完整用户指针内容由硬件导出 & \nope{NO} & 硬件 pointer snapshot 是有界前缀，完整字符串主要靠受信 companion 对齐。 \\
    \rowcolor{black!4}
    board trace 直接提供源码行号 & \nope{NO} & 当前板上 trace 是函数级归属；source-line 是 sidecar 证据。 \\
    全系统连续语义恢复 & \nope{NO} & 目前结论限定在捕获窗口和样例范围内。 \\
    \bottomrule
  \end{tabularx}
  \end{center}
  \vskip0.08cm
  \keybox{主动讲边界不是削弱结果，而是让当前进展更可信。}
\end{frame}

% 27
\begin{frame}{两个最容易被问到的边界}
  \vskip0.1cm
  \begin{columns}[T,onlytextwidth]
    \column{0.48\textwidth}
    \begin{block}{pointer 参数}
      现在可以说：有界前缀信息和 companion 对齐可以支撑当前样例语义。不能说：硬件已经完整导出用户字符串或结构体。
    \end{block}
    \column{0.48\textwidth}
    \begin{block}{源码行号}
      现在可以说：函数级归属可用。不能说：板上 trace 已经是 DWARF 行号 trace。
    \end{block}
  \end{columns}
  \vskip0.25cm
  \keybox{如果老师问“这到底证明了什么”，建议用“窗口内记录、行为解释、边界清楚”三句话回答。}
\end{frame}

\section{下一阶段}
\sectiondivider{五}{下一阶段}

% 29
\begin{frame}{下一阶段不建议盲目加样例}
  \vskip0.1cm
  样例数量继续增加当然有价值，但现在更重要的是把论文主线收窄。需要回答的是：硬件辅助 trace 相比纯软件方法有什么价值；哪些语义必须由硬件补；哪些语义可以由受信 helper 补。

  \vskip0.25cm
  \begin{center}
  \renewcommand{\arraystretch}{1.18}
  \begin{tabularx}{0.92\linewidth}{@{}L{3.1cm}Y@{}}
    \toprule
    \textbf{硬件价值} & 旁路记录、低可见度、DROP 可见、与软件对照方法区分。 \\
    \rowcolor{black!4}
    \textbf{语义价值} & 在有证据的范围内恢复路径、进程、动态映射和反调试行为。 \\
    \textbf{论文价值} & 明确贡献边界，不把安全替代样例包装成真实恶意检测。 \\
    \bottomrule
  \end{tabularx}
  \end{center}
\end{frame}

% 30
\begin{frame}{近期任务安排}
  \vskip0.05cm
  \begin{center}\scriptsize
  \renewcommand{\arraystretch}{1.18}
  \begin{tabularx}{0.96\linewidth}{@{}L{3.0cm}Y L{2.2cm}@{}}
    \toprule
    \textbf{任务} & \textbf{交付物} & \textbf{判断标准} \\
    \midrule
    hardware pointer 路线 & 明确哪些参数能由 \texttt{ARG\_MEM} 有界补充。 & 路线可决定 \\
    \rowcolor{black!4}
    Linux workload 固定 & open/exec/mmap/process 的可复查脚本。 & 语义可复查 \\
    对照实验整理 & strace、QEMU、helper、软件插桩的同口径指标。 & 表格可引用 \\
    \rowcolor{black!4}
    论文边界表 & allowed claims / non-claims 写成论文式表格。 & 可直接审阅 \\
    复现实验包 & 命令、hash、摘要和检查脚本。 & 他人可复跑 \\
    \bottomrule
  \end{tabularx}
  \end{center}
  \vskip0.08cm
  \keybox{优先补“论文里一定会被问”的证据缺口，再扩展样例规模。}
\end{frame}

% 31
\begin{frame}{需要老师帮忙定的三个方向}
  \vskip0.15cm
  \begin{enumerate}\setlength\itemsep{0.65em}
    \item \textbf{论文定位}：主打 RISC-V 行为追踪，还是主打恶意行为分析应用。
    \item \textbf{pointer 优先级}：继续推进硬件补全，还是把 helper 路线作为可信语义补充。
    \item \textbf{近期目标}：先形成技术报告/小论文，还是继续补强到更完整投稿。
  \end{enumerate}
  \vskip0.25cm
  \keybox{\textbf{当前状态：}阶段汇报已经有充分数据支撑；下一步需要决定论文主线和补实验顺序。}
\end{frame}

\begin{frame}[plain]
  \begin{tikzpicture}[remember picture,overlay]
    \fill[accentcolor!7] (current page.north west) rectangle (current page.south east);
    \fill[accentcolor] (current page.north west) rectangle ([yshift=-0.16cm]current page.north east);
    \fill[accentcolor] ([yshift=0.16cm]current page.south west) rectangle (current page.south east);
    \node[text=accentcolor,font=\sffamily\bfseries\fontsize{26pt}{34pt}\selectfont]
      at ([yshift=0.35cm]current page.center) {讨论与建议};
    \node[text=textgray,font=\sffamily\fontsize{10pt}{14pt}\selectfont,align=center]
      at ([yshift=-1.0cm]current page.center) {RV-MalTrace 阶段性进展汇报 \;|\; 2026年6月12日};
  \end{tikzpicture}
\end{frame}

\end{document}
